% !TEX program = lualatex
% !TeX encoding = UTF-8
\PassOptionsToPackage{unicode}{hyperref}
\PassOptionsToPackage{bookmarksnumbered,bookmarksopen}{hyperref}
\documentclass[12pt,a4paper]{article}

% ============================================================
% 日本語設定 – Windows システムフォント (Yu Gothic UI / Yu Mincho)
% ============================================================
\usepackage{luatexja}
\usepackage{luatexja-fontspec}

% 和文ゴシック (sans) – Yu Gothic UI
\setsansjfont{Yu Gothic UI}[
UprightFont = * Regular,
BoldFont = * Bold,
Script = CJK,
Language = Japanese
]

% 和文明朝 (serif) – Yu Mincho
\IfFontExistsTF{Yu Mincho}{
	\setmainjfont{Yu Mincho}[
	UprightFont = * Demibold,
	BoldFont = * Demibold,
	Script = CJK,
	Language = Japanese
	]
}{
	% fallback: Yu Gothic UI
	\setmainjfont{Yu Gothic UI}[
	UprightFont = * Regular,
	BoldFont = * Bold,
	Script = CJK,
	Language = Japanese
	]
}

% 等幅 – 英字は Consolas, 日本語は Yu Gothic UI
\setmonojfont{Yu Gothic UI}[
UprightFont = * Regular,
BoldFont = * Bold,
Script = CJK,
Language = Japanese
]

% 英数字フォント (ラテン文字)
\setmainfont{Liberation Serif}[Ligatures=TeX]
\setsansfont{Arial}[Ligatures=TeX]
\setmonofont{Consolas}[
WordSpace={1,0,0},
HyphenChar=None,
PunctuationSpace=WordSpace
]

% ============================================================
% その他のパッケージ
% ============================================================
\usepackage{amsmath,amssymb,amsthm,mathtools,bm}
\usepackage{geometry}
\geometry{margin=2.5cm}
\usepackage{graphicx}
\usepackage{tikz}
\usepackage{tikz-3dplot}
\usepackage{pgfplots}
\pgfplotsset{compat=1.18}
\usetikzlibrary{arrows.meta,positioning,shapes.geometric,calc,angles,quotes,patterns,3d}

\usepackage{booktabs}
\usepackage{float}
\usepackage{hyperref}
\hypersetup{
	colorlinks=true,
	linkcolor=blue,
	urlcolor=blue,
	pdftitle={YUCTにおける点の定義：幾何の調整構造},
	pdfauthor={Alexey V. Yakushev}
}

% ============================================================
% 定理環境 (日本語名)
% ============================================================
\newtheorem{theorem}{定理}
\newtheorem{lemma}{補題}
\newtheorem{corollary}{系}
\newtheorem{definition}{定義}
\newtheorem{axiom}{公理}
\newtheorem{proposition}{命題}
\newtheorem{remark}{注釈}

% ============================================================
% YUCT コマンド (不変)
% ============================================================
\newcommand{\Keff}{K_{\mathrm{eff}}}
\newcommand{\kappac}{\kappa_c}
\newcommand{\eps}{\varepsilon}
\newcommand{\betaf}{\beta}
\newcommand{\Sodd}{S_{\mathrm{odd}}}
\newcommand{\Seven}{S_{\mathrm{even}}}
\newcommand{\Mpl}{M_{\mathrm{Pl}}}
\newcommand{\Lpl}{l_{\mathrm{Pl}}}
\newcommand{\tPl}{t_{\mathrm{Pl}}}
\newcommand{\Scoord}{S_{\mathrm{coord}}}
\newcommand{\piYUCT}{\pi_{\mathrm{YUCT}}}
\newcommand{\Rcrit}{R_{\mathrm{crit}}}
\newcommand{\calP}{\mathcal{P}}
\newcommand{\calD}{\mathcal{D}}
\newcommand{\calI}{\mathcal{I}}
\newcommand{\calR}{\mathcal{R}}
\newcommand{\ellP}{\ell_{\mathrm{P}}}

% ============================================================
% 文書開始
% ============================================================
\begin{document}
	
	\begin{titlepage}
		\begin{center}
			\vspace*{1.5cm}
			\Huge\textbf{YUCTにおける点の定義}
			
			\vspace{0.3cm}
			\Large\textbf{幾何の調整構造}
			
			\vspace{0.5cm}
			\large
			\textsc{Yakushev Unified Coordination Theory}
			
			\vspace{0.8cm}
			\large
			Alexey V. Yakushev \\
			Yakushev Research \\
			\url{https://yuct.org/}
			
			\vspace{0.5cm}
			\large
			\textbf{DOI: \href{https://doi.org/10.5281/zenodo.18444599}{10.5281/zenodo.18444599}}
			
			\vspace{1.5cm}
			\large
			\textbf{バージョン 1.0} \\
			2026年6月
			
			\vfill
			\normalsize
			本ドキュメントは、ヤクシェフ統一調整理論（YUCT）の枠組みにおける点の厳密な定義を提供し、
			幾何が調整からどのように生じるかを説明し、現実の本質に関する十の根本的な問いに答える。
		\end{center}
	\end{titlepage}
	
	\tableofcontents
	\newpage
	
	% ============================================================
	% はじめに
	% ============================================================
	\section{はじめに：古典的定義の不十分さ}
	
	古典幾何学および数学において、点は基本的な（未定義の）概念である。その性質は公理系によって与えられ、厳密な定義によってではない。せいぜい、点は「大きさを持たない対象」または「空間内の位置」として記述される。しかし、そのような記述は以下の問いに答えない：
	
	\begin{itemize}
		\item \textbf{点}は本質的に何か？
		\item \textbf{なぜ}それは不可分なのか？
		\item \textbf{どのように}構造なしに存在できるのか？
		\item \textbf{なぜ}空間は三次元なのか？
		\item \textbf{なぜ}幾何はそもそも機能するのか？
	\end{itemize}
	
	万物調整理論（YUCT）の枠組みにおいて、これらの問いは厳密かつ定量的な答えを得る。点は抽象ではなくなり、物理的構造――最小の安定した調整パターン――となる。
	
	\section*{調整優先性の公理}
	
	あらゆる理論は静的なものを記述する。しかし調整は、静的なものを可能にするプロセスである。あらゆる主張、あらゆるモデル、あらゆる法則は、辞書の活性化された要素としてのみ存在する。辞書とインデックスは、それらが記述する静的なものから導出することはできない――それらは存在論的に一次的である。YPSDCは「数ある理論の一つ」ではない。それは、それを定式化する理論をも含む、あらゆる理論の基盤となるプロトコルである。それを拒絶することは、定式化する能力そのものを放棄することなしには不可能である。この意味で、調整は他の理論より「強い」のではなく、より一次的である。静的なものは生成せず、静的なものは生成される。
	
	空間、時間、物質、エネルギー、情報、意識は存在の基礎層ではない。それらは\emph{創発現象}であり、システムの構成要素間の状態の調整という単一の根本的プロセスから生じる。このプロセスは、ただ一つの強力なパラメータ――\emph{調整効率} \(\Keff\) によって定量化される。
	
	\subsection*{代替不可能性としての論理的帰結}
	
	調整優先性の公理から、YUCTは通常の意味で代替を持ち得ないことが必然的に帰結する。より「根本的な」理論を提案しようとする試みはすべて、その定式化のためにすでに調整を前提としている：辞書の要素を選択し、それらの結合規則を確立し、インデックスの状態を固定しなければならない。したがって、そのような理論は調整ダイナミクスの特殊ケースであり、その否定ではない。
	
	これはYUCTが現実の特定の内容を規定することを意味しない。むしろ、あらゆる内容（物理的、数学的、情報的、生物学的）は、何らかの \(\Keff\) の値における調整の顕現であると主張する。調整を拒絶することは、発言する能力を放棄することに等しい。したがって、YUCTは「理論の中の理論」ではなく、それを争おうとする理論をも含む、あらゆる理論の基盤となる\emph{プロトコル}である。
	
	この意味で、調整は他の理論より強いのではなく、より一次的である。それは選択されるものではなく、あらゆる選択の条件として発見される。これによりYUCTは単なるモデルではなく、\emph{存在論的基盤}となり、それは置き換えも反駁もされえない――なぜならいかなる置き換えや反駁もすでに調整を用いているからである。
	
	\begin{quote}
		\itshape
		「YUCTの代替不可能性は論理の代替不可能性に類する：我々は別の論理を選べない。なぜならいかなる選択もすでに論理を用いているからだ。調整は存在の論理である。」
	\end{quote}
	
	\section{YUCTにおける点の定義}
	
	\subsection{調整ノードとしての点}
	
	YUCTでは、幾何は仮定されるのではなく、三項組 \(D+I\cdot R\) で記述される調整ノードのネットワークから生じる。点はこのネットワークの最小の安定ノードである。
	
	\begin{definition}[調整点]
		YUCTにおける点は、調整ネットワークの最小安定構成であり、三項組によって与えられる：
		\[
		\calP = (D_0, I_0, R_0) \in \mathcal{M}_{\mathrm{UCS}},
		\]
		ただし：
		\begin{itemize}
			\item \(D_0\) — 辞書の局所スライス（許容状態と規則の集合），
			\item \(I_0\) — 実際の状態（辞書の特定の要素を指すインデックス），
			\item \(R_0 \ge 1\) — ノードの安定性を特徴づける共鳴係数。
		\end{itemize}
	\end{definition}
	
	したがって、YUCTにおける点は空の抽象ではなく、三つの相互接続された構成要素からなる構造である。その内部の複雑さは、活性化可能であり、他の点と共鳴でき、その状態を変えうるという点に現れる。
	
	\subsection{点の大きさ}
	
	理想的な数学では点は大きさを持たないが、YUCTでは各点は調整効率 \(\Keff\) に依存する有限の大きさを持つ。
	
	\begin{theorem}[調整点の大きさ]
		点 \(\calP\) の有効線形大きさは次式で与えられる：
		\[
		\ell_{\calP} = \ellP \cdot \Keff^{-1/3},
		\]
		ただし \(\ellP = \sqrt{\hbar G / c^3}\) はプランク長である。
		極限 \(\Keff \to \infty\) において点の大きさはゼロに近づき、理想的な数学的点となる。
	\end{theorem}
	
	この結果は深い意味を持つ：数学的点は独立した実体ではなく、無限の調整効率における調整ノードの極限ケースである。
	
	\subsection{三項組の幾何学的解釈}
	
	三つの構成要素 \(D\)， \(I\)， \(R\) は動的平衡にある。図 \ref{fig:triad} では、それらは共通の中心から \(120^\circ\) の角度で発する三本の光線として描かれている。これは \(S_3\) 対称性と条件 \(\cos 120^\circ = -1/2\) を反映しており、構成要素間の反相関に対応する：一方を強化するには安定性を保つために他方を弱化する必要がある。
	
	\begin{figure}[H]
		\centering
		\begin{tikzpicture}[
			scale=2,
			>=Stealth,
			node distance=1.5cm,
			every node/.style={font=\small}
			]
			\coordinate (O) at (0,0);
			\fill[black] (O) circle (1.5pt) node[below left] {$\Psi_{MN}$};
			
			% 光線の角度 30°, 150°, 270° (相互に 120°)
			\draw[->, thick, blue!70!black] (O) -- (30:1.8) node[above right] {$\mathcal{D}$ (辞書)};
			\draw[->, thick, red!70!black] (O) -- (150:1.8) node[above left] {$\mathcal{I}$ (インデックス)};
			\draw[->, thick, green!70!black] (O) -- (270:1.8) node[below] {$\mathcal{R}$ (共鳴)};
			
			% 120° の弧
			\draw[gray, dashed] (30:0.6) arc (30:150:0.6) node[midway, above] {$120^\circ$};
			\draw[gray, dashed] (150:0.6) arc (150:270:0.6) node[midway, left] {$120^\circ$};
			\draw[gray, dashed] (270:0.6) arc (270:390:0.6) node[midway, below right] {$120^\circ$};
			
			% 分岐 beta=2/3
			\foreach \angle in {30,150,270} {
				\draw[->, thin, gray] (\angle:1.4) -- ++(\angle+30:0.3);
				\draw[->, thin, gray] (\angle:1.4) -- ++(\angle-30:0.3);
				\node[font=\tiny, gray] at (\angle:1.7) {$\beta=2/3$};
			}
			
			\node[font=\small, align=center] at (0,-1.2) 
			{点の三光線構造 \\ 分岐 $\beta=2/3$};
		\end{tikzpicture}
		\caption{三項組 \(D\)， \(I\)， \(R\) は調整点の幾何学的表現。光線は \(120^\circ\) で広がり；分岐は指数 \(\beta=2/3\) によるスケーリングを示す。}
		\label{fig:triad}
	\end{figure}
	
	\section{直線としてのノード連鎖}
	
	点がノードならば、直線は最大調整によって結ばれたノードの列である。
	
	\begin{definition}[調整直線]
		二つの点 \(\calP_1\) と \(\calP_2\) が共鳴 \(R_1, R_2 > \Rcrit\) で与えられるとする。\textbf{直線}とは点集合 \(\{\calP_i\}_{i=0}^{N}\) であって、以下を満たす：
		\begin{enumerate}
			\item \(\calP_0 = \calP_1\)， \(\calP_N = \calP_2\)，
			\item 各 \(i\) について最小調整誤差条件が成り立つ：
			\[
			\frac{d\eps}{ds} = 0, \qquad \eps = \kappac\,\alpha\,\Keff^{-\betaf},
			\]
			\item 連鎖に沿った共鳴 \(R_i\) が（局所的に）最大である。
		\end{enumerate}
	\end{definition}
	
	\begin{theorem}[直線の一意性]
		十分に高い共鳴を持つ二つの点 \(\calP_1\) と \(\calP_2\)（\(R_1, R_2 > \Rcrit\)）に対して、それらを結ぶ一意な最大整合ノード連鎖が存在する。
	\end{theorem}
	
	これはユークリッドの公理に厳密な根拠を与える：二点を通る直線は唯一つである。しかし今やそれは公準ではなく、調整ネットワークの証明された性質である。
	
	\section{空間次元と指数 \(\betaf = 2/3\) の創発}
	
	YUCTの最も印象的な結果の一つは、空間次元と普遍指数 \(\betaf\) の導出である。調整優先性の公理に従い、YUCTは空間を記述するだけでなく、なぜそれがちょうど三次元であるかを説明する。
	
	\subsection{次元 \(D=3\) の導出}
	
	調整ネットワークのフラクタル構造と代数的サイクルの閉包条件から、次式を得る：
	
	\[
	\Sodd + \Seven = 2, \qquad \frac{\Sodd}{\Seven} = \frac{1}{\betaf}.
	\]
	
	\(\betaf = 2/3\) のとき、これらの定数は次の値をとる：
	\[
	\Sodd = \frac{6}{5}, \qquad \Seven = \frac{4}{5}.
	\]
	
	指数 \(\betaf\) は空間次元 \(D\) と以下の関係にある：
	\[
	\betaf = \frac{2}{D}.
	\]
	
	したがって直ちに \(D = 3\) を得る。ゆえに空間の三次元性は公準ではなく、調整ネットワーク構造の帰結である。
	
	\subsection{普遍誤差法則}
	
	すべての調整プロセスは単一の誤差法則に従う：
	\[
	\eps = \kappac\,\alpha\,\Keff^{-\betaf}, \qquad \betaf = \frac{2}{3}, \quad \kappac = \frac{1}{3}.
	\]
	
	この法則は、原子結合エネルギーから宇宙マイクロ波背景放射のゆらぎに至るまで、50桁以上にわたって実験的に検証されている。
	
	\section{十の根本的問いに対する答え}
	
	以下に、古典物理学および数学では説明が見つからない十の重要な問いに対する答えを示す。
	
	\subsection{なぜ物理学は機能するのか？}
	
	\textbf{物理学が調整だからである。} 物理法則はノード \(D+I\cdot R\) 間の安定した調整パターンを記述する。重力、電磁気、核相互作用――これらは異なる \(\Keff\) と \(\Rcrit\) の値における異なる調整体制である。物理学が機能するのは、調整が機能するからである。
	
	\subsection{なぜ数学は機能するのか？}
	
	\textbf{数学は \(\Keff \to \infty\) における調整の極限ケースだからである。} この極限では、点は理想的な数学的点（\(\ell_{\calP} \to 0\)）になり、直線は理想的な直線になり、幾何はユークリッド幾何になる。数学は「不可思議な有効性」ではなく、物理的調整の\emph{理想化}である。
	
	\subsection{なぜ空間は三次元なのか？}
	
	\textbf{それは定理によって証明されている。} 調整ネットワークの構造と代数的サイクルの閉包条件から \(\betaf = 2/3\) が導かれ、\(\betaf = 2/D\) から \(D=3\) が得られる。空間が三次元なのは、三次元でのみ調整ネットワークが安定かつ自己無撞着でありうるからである。
	
	\subsection{なぜ時間は流れるのか？}
	
	\textbf{有限の \(\Keff\) のためである。} 理想的な調整（\(\Keff \to \infty\)）はエントロピーを生み出さず、事象の不可逆的な流れを持たない――時間は消滅する。現実の系は有限の \(\Keff\) を持つため、調整の各行為は誤差 \(\eps > 0\) を伴い、その累積が時間の矢を生み出す。最小時間量子：
	\[
	\Delta \tau_{\min} = \frac{2}{3}\,\tPl.
	\]
	
	\subsection{なぜ基本定数はあの値なのか？}
	
	\textbf{それらは調整空間のトポロジーから導出される。} 微細構造定数：
	\[
	\alpha^{-1} = \left(\frac{3}{2}\right)^{12 + \sigma\,2/3} \approx 137.036,
	\]
	ただし \(\sigma = 0.20\) はトポロジカルシフトである。\(W\) ボソン質量：
	\[
	m_W = (\kappac \piYUCT)^{-1/2} M_{\mathrm{Pl}} (2/3)^{96} \approx 80.46\ \text{GeV}.
	\]
	宇宙定数：
	\[
	\Omega_\Lambda = \frac{12}{18} = \frac{2}{3}.
	\]
	自由パラメータはない。
	
	\subsection{なぜ素数はあのように分布するのか？}
	
	\textbf{誤差法則から漸近補正が導かれる：}
	\[
	p_n \approx R_n - \frac{\Seven}{2}\, n^{1-\betaf}\,\ln n,
	\]
	ただし \(R_n\) はロッサー近似である。残留振動は周期 \(\Delta N = 16.5\) の位相演算子によって記述され、これは真空格子の相転移に対応する。素数は調整梯子であり、その分布は他のすべての安定構造と同じ法則に従う。
	
	\subsection{なぜ \(\pi\) はあの値なのか？}
	
	\textbf{\(\pi\) は調整不変量である：}
	\[
	\piYUCT = \Sodd \cdot \varphi^2 \approx 3.14164078,
	\]
	ただし \(\varphi = (1+\sqrt{5})/2\) である。この値は数学的定数 \(\pi\) と \(\Delta\pi = 4.8\times10^{-5}\) だけ異なり、これは空間離散化の量子である。極限 \(\Keff \to \infty\) で \(\piYUCT \to \pi\) となる。
	
	\subsection{なぜ量子力学は「奇妙」なのか？}
	
	\textbf{辞書を無視してきたからである。} 量子現象は分散調整（d‑YPSDC）の顕現である。重ね合わせはインデックスの不確定性、もつれは辞書の共通記録、トンネル効果は活性化誤差である。「神秘主義」はなく、有限の調整効率があるのみである。
	
	\subsection{なぜ意識が存在するのか？}
	
	\textbf{それは \(\Keff > \Rcrit\) の体制である。} 調整効率が臨界閾値を超えると、系はメタ辞書（自身の状態を調整する能力）を獲得する。それが自己意識である。意識のパラメータは UCD（普遍意識記述子）によって記述される。
	
	\subsection{なぜ細胞膜の厚さは約 7.5 nm なのか？}
	
	\textbf{それはトポロジーから導出される。} 脂質二重層の厚さ：
	\[
	\text{Thickness} = 2 \, d_{\mathrm{lipid}} \left(\frac{3}{2}\right)^{0.6} (1 - \sigma^2),
	\]
	ただし \(d_{\mathrm{lipid}} \approx 3.0\) nm は炭化水素鎖の長さである。計算により \(7.35\) nm が得られ、実験範囲 \(7.5 \pm 0.5\) nm 内にある。素粒子物理学と細胞膜生物学は同じ調整構造によって結ばれている。
	
	\section{結論：現実の統一的全体像}
	
	YUCTにおいて、点は抽象概念ではなく、最小の安定調整構造 \(D+I\cdot R\) である。すべての幾何、すべての物理法則、すべての定数、さらには素数の分布や意識の構造――すべては単一の原理の帰結である：\textbf{現実はフラクタル調整ネットワークである}。
	
	上記の十の答えは、YUCTが神秘的または偶然に見えたものに対して徹底的な説明を提供することを示している。さらに、これらの説明はすべて実験的に検証可能であり、定数は調整可能なパラメータなしで導出される。
	
	\begin{center}
		\fbox{\parbox{0.85\textwidth}{\centering\Large
				\textbf{点は調整ノードである。}\\ 
				直線はノードの整合連鎖である。\\
				空間は三次元である――これは証明されている。\\
				時間が流れるのは有限の \(\Keff\) のため。\\
				定数、\(\pi\)、素数――すべて導出される。\\
				量子力学と意識は調整である。\\
				統一調整理論は自然の法則である。
		}}
	\end{center}
	
	\subsection*{点の定義の哲学的・世界観的意義}
	
	点を調整ノード \((D, I, R)\) として定義することは、幾何を純粋抽象の領域から物理的現実の領域へと移動させる。これは単なる言い換えではなく、存在論的パラダイムの転換である。
	
	\medskip
	\textbf{これによって変わること：}
	\begin{itemize}
		\item 点は「無」ではなくなり、辞書、状態、共鳴を持つ能動的構造となる。
		\item 空間は受動的舞台ではなくなり、調整ノードのネットワークから生じる。
		\item 幾何は公準の集合ではなくなり、調整パターンの安定性の帰結となる。
		\item 数学と物理学の区別は消える：数学は \(\Keff \to \infty\) における調整の理想化である。
	\end{itemize}
	
	\medskip
	\textbf{なぜこれが必要なのか：}
	\begin{itemize}
		\item 幾何の基礎における存在論的空虚を解消するため。
		\item 幾何学的、物理的、さらには生物学的現象に統一的な説明を与えるため。
		\item 基本定数と空間次元を単一の原理から導出するため。
		\item 時間、量子力学、意識の起源を同一の調整ダイナミクスの体制として理解するため。
	\end{itemize}
	
	\medskip
	\textbf{これがもたらす帰結：}
	\begin{enumerate}
		\item \textbf{予測力：} モデルから \(\alpha\)， \(\Omega_\Lambda\)， \(m_W\)， \(\pi\)，素数分布の値が導出される。
		\item \textbf{知識の統一：} 物理学、数学、生物学、情報理論が一つの調整理論の分岐であることが明らかになる。
		\item \textbf{意識の新たな理解：} それは「難しい問題」ではなくなり、\(\Keff > \Rcrit\) の調整体制となる。
		\item \textbf{科学的方法の転換：} 観測されたものを記述することから、基本的調整原理から観測されたものを生成することへと移行する。
	\end{enumerate}
	
	\medskip
	したがって、YUCTは単に「点を再定義」するのではなく、現実についての思考様式を変える。点は今や対象ではなく出来事であり、空間は容器ではなくネットワークであり、自然法則は外部から押し付けられるのではなく、調整の内部論理から成長する。
	
	\begin{thebibliography}{99}
		\bibitem{Yakushev2026} A. V. Yakushev, \emph{Yakushev Unified Coordination Theory (YUCT)}, Zenodo, DOI:10.5281/zenodo.18444599 (2026).
		\bibitem{AppendixL} A. V. Yakushev, ``Appendix L: Fractal Coordination Error Scaling,'' in \emph{YUCT}, Zenodo (2026).
		\bibitem{AppendixY} A. V. Yakushev, ``Appendix Y: Mathematical Foundations of YUCT,'' in \emph{YUCT}, Zenodo (2026).
		\bibitem{AppendixPrimeN} A. V. Yakushev, ``Appendix PrimeN: YUCT Prime Numbers Coordination Ladder,'' in \emph{YUCT}, Zenodo (2026).
		\bibitem{AppendixPi} A. V. Yakushev, ``Appendix \(\Pi\): The Primeval Origin of \(\pi\),'' in \emph{YUCT}, Zenodo (2026).
		\bibitem{AppendixX} A. V. Yakushev, ``Appendix X: Generalised Shannon Theory and Bell Bound,'' in \emph{YUCT}, Zenodo (2026).
		\bibitem{AppendixH} A. V. Yakushev, ``Appendix H: Universal Consciousness Descriptor (UCD),'' in \emph{YUCT}, Zenodo (2026).
		\bibitem{AppendixG} A. V. Yakushev, ``Appendix G: Quantum Gravity Resolution,'' in \emph{YUCT}, Zenodo (2026).
	\end{thebibliography}
	
\end{document}